\section{Introduction}
Learned world models support robot control by predicting the consequences of candidate actions. Visual models make this approach possible directly from images, using a learned representation in which a planner can compare alternative futures~\cite{zhou2025dinowm}. Deployment, however, changes the relationship between action and outcome. A familiar object slides differently on another table, a container carries a different load, or a replacement tool changes the robot's effective dynamics. Online adaptation addresses these discrepancies by updating a model from the errors observed during execution. Recent approaches update world-model parameters or lightweight residual modules using transitions collected while planning~\cite{wang2026adajepa,soni2026sandwich}.

This strategy has a timing limitation. Consider a robot striking a puck toward a target. The appropriate strike depends on the surface and on the object's response to impact. If the robot assumes too much friction, the puck may overshoot the target or leave the table. The prediction error reveals the mismatch, but only after the strike has been committed. Correcting the model at that point can improve the next attempt; it cannot change the impulse already delivered. Similar timing constraints arise in throwing, in a lift whose initial force can cause slipping or excessive acceleration, and in pushing sufficiently close to an edge that the initial motion is consequential. These are \emph{single-contact tasks} in the sense relevant here: an outcome component is determined by an action committed before feedback from that interaction can affect the decision.

For an adaptation rule driven exclusively by current-episode prediction errors, the limitation is immediate. Before any eligible error feedback has arrived, the predictor remains at its episode initialization. If that initialization is the frozen model, its first-contact prediction is exactly the frozen prediction. This observation does not depend on the optimizer, the number of parameters being adapted, or how quickly adaptation proceeds once data arrive. It identifies an information boundary: improving the consequential first prediction requires useful information available before the error it would otherwise generate.

Many deployments provide such information through recurrence. A robot may repeatedly encounter the same surfaces, payload classes, containers, and tools. Treating each return as a fresh estimation problem discards knowledge that could already be useful. In this setting, adaptation includes a second question alongside parameter estimation: \emph{which previously encountered physical context is active now?} Maintaining separate dynamics memories makes this question operational. A robot can preserve what it learned about one condition while operating in another, then retrieve that knowledge when the earlier condition returns. Persistent mixtures of dynamics models already establish a precedent for this form of reuse~\cite{nagabandi2019mole}; the first-contact setting asks when such knowledge can influence an action before new dynamics evidence is available.

Contextual cues provide one possible answer. A surface may look like felt or acrylic, a container may identify a recurring payload class, and a visible tool may indicate which actuator response to expect. These cues are imperfect evidence. Identical appearances can conceal different physics, and familiar physics can acquire a new appearance. We therefore use cues to infer a distribution over robot-discovered dynamics memories. The robot learns what a cue predicts through interaction, and subsequent dynamics evidence can confirm or overturn the resulting belief. The relevant operation is retrieval of learned physical experience, rather than an assumption that appearance uniquely determines physics. A truly novel condition with no transferable structure remains uncertain before it is encountered.

Recurrence also has compositional structure. A robot may have experienced surface B with a light load and a heavy load on surface A, yet never encountered surface B with the heavy load. Independent memories indexed only by the complete condition cannot directly reuse these experiences as a new combination. A shared prior over factor effects can transfer their contributions, provided the effects are sufficiently separable and previous experience identifies the required relations. This leads to a stronger test than rapid adaptation after a change: prediction at the first contact of the first-ever occurrence of a held-out combination. Independent interventions do not guarantee additive dynamics, so the prior also represents combination-specific interactions and the uncertainty they leave unresolved.

We introduce \method{} (Context-Indexed Residual Adapters), a persistent library of compact probabilistic corrections around a frozen visual world model. Before contact, a context posterior combines deployment history with the current cue to retrieve a distribution over corrections. After contact, dynamics likelihoods revise this belief, and sustained evidence of unexplained dynamics supports a new memory. A shared Gaussian prior organizes complete-context memories by the physical factors they have in common. This gives an unseen combination a predictive distribution before its first interaction, including uncertainty about its unobserved interaction component.

Learning those memories introduces a second role for context inference. A model that stores an error according to how well that same error matches a memory can systematically select its observation noise. \method{} therefore weights each update by the context probabilities available before observing the current residual. We derive a uniform memory-error bound for this predictable update. Its terms distinguish observation noise, mismatch in the transferred prior, and contamination by incorrectly assigned contexts. The context filter's log loss controls the accumulated contamination and also bounds first-contact selection error. Under Lipschitz dynamics and costs, these errors translate into a gap relative to the best candidate committed action. The analysis thus links the quality of a contextual repertoire to the decisions it can support.

Contextual uncertainty also determines whether additional interaction is useful. In a separate probe-enabled protocol, the robot can gather evidence before committing to the strike. A probe is selected by its expected effect on the decision, accounting for its cost. A value-of-information result shows why identifying the physical condition need not help: if all plausible models share an optimal action, informative observations have zero decision value. The same posterior supports a safety filter with per-memory calibration, so a familiar condition can recover a previous uncertainty margin. A long-run bound concerns failures of the latent barrier condition under explicit boundedness and enforcement assumptions; it does not establish physical safety from context recognition alone.

The central empirical question is whether this architecture changes the mechanism of adaptation. A faster loss decrease alone cannot distinguish retrieving an existing model from fitting a new one. A counterfactual replay decomposition separates changes in memory selection from changes in memory parameters and transferred prior centres. A complementary intervention freezes learned memories during later returns. Controlled changes to cue informativeness, recurrence, factor combinations, and intervening experience isolate when persistent context inference helps, while calibration and probing experiments test the additional predictions of the analysis.

The resulting perspective changes the unit of adaptation. Deployment becomes inference over a growing repertoire of physical contexts, with parameter learning used to acquire or refine the memories in that repertoire. Recurrence then allows previously learned dynamics to influence the first action of a new episode, while composition allows familiar physical attributes to contribute predictions in unfamiliar combinations. This paper develops four contributions:
\begin{itemize}[leftmargin=*]
  \item \textbf{First-contact adaptation and control.} We formalize the timing limitation of error-driven episode resets and connect contextual prediction accuracy to the cost of the committed action.
  \item \textbf{Persistent compositional memory.} We introduce cue-guided residual retrieval with a shared factor prior and characterize when observed combinations identify a new additive combination, including the uncertainty floor left by independent interactions.
  \item \textbf{Learning and acting under context uncertainty.} We bound memory error under predictable context-weighted updates, characterize decision-relevant probing, and establish a conditional long-run bound for memory-specific latent safety calibration.
  \item \textbf{Mechanisms of adaptation.} We separate recall from parameter learning and define controlled evaluations of first contact, recurrence, composition, interference, calibration, and information gathering.
\end{itemize}
